\documentclass[12pt]{article}
\usepackage[utf8]{inputenc}
\usepackage[a4paper, margin=2cm]{geometry}
\usepackage{amsthm}
\usepackage{amsmath, amssymb}
\usepackage{circuitikz}
\usepackage{graphicx}
\usepackage{karnaugh-map}
\usepackage{newtxtext, newtxmath}
\usepackage{tcolorbox}
\usepackage{tikz}
\usepackage{titlesec}
\usepackage{xcolor}
\usepackage{xeCJK}
\usepackage{listings}
\definecolor{lightblue}{RGB}{135, 206, 250}
\tcbuselibrary{skins,theorems}

\titleformat{\section}[block]{\Large\bfseries\centering}{}{0pt}{}
\titleformat{\subsection}[block]{\bfseries}{\thesubsection}{2pt}{}

\newtcbtheorem[no counter]{proposition}{\large\hspace{-8pt}Proposition}{
  enhanced,
  colback=white,
  colframe=lightblue,
  left=12pt,
  right=12pt,
  top=24pt,
  bottom=24pt,
  fonttitle=\bfseries
}{prop}
\newtcbtheorem[]{question}{\large\hspace{-8pt}Question :}{
  enhanced,
  colback=white,
  colframe=black,
  left=12pt,
  right=12pt,
  top=24pt,
  bottom=24pt,
  fonttitle=\bfseries
}{qst}
\newenvironment{Answer}{
  \vspace{12pt}
  {\large\textbf{\textit{\hspace{-12pt}Answer:}}}\par
  \vspace{8pt}
} {\vspace{24pt}}
\newenvironment{Proof}{
  \vspace{12pt}
  {\large\textbf{\textit{\hspace{-12pt}Proof:}}}\par
  \vspace{8pt}
} {\vspace{24pt}}

\lstset{
    language=Python,
    basicstyle=\ttfamily\footnotesize,
    keywordstyle=\color{blue}\bfseries,
    commentstyle=\color{gray},
    stringstyle=\color{red},
    showstringspaces=false,
    numberstyle=\tiny\color{gray},
    stepnumber=1,
    numbersep=10pt,
    frame=single,
    breaklines=true,
    captionpos=b
}

\title{\textbf{Algorithm Design Homework 2}}
\date{June 18, 2024}

\begin{document}
\maketitle

\begin{question}{}{}
给定一个0-1背包问题weight w = \{2,1,3,2\}, value v = \{12,10,20,15\}, capacity c = 5。
如何使动态规划算法在O(c)空间内得到最优解(即n元0-1向量) ?
\begin{itemize}
  \item (1) 给出得到最优解的计算过程。
  \item (2) 给出该过程所需的空间复杂度。
  \item (3) 给出该过程的时间复杂度。
\end{itemize}
\end{question}
\begin{Answer}
\subsection*{计算过程}
\noindent

首先，定义动态规划数组 dp，其中 dp[j] 表示容量为 j 时可以获得的最大价值。

初始化 dp 数组：

$$dp=[0]*(c + 1)$$

我们遍历每个物品，并更新\texttt{dp}数组，从背包的最大容量开始往前更新。

物品权重和价值分别为：
\begin{itemize}
  \item weight = [2, 1, 3, 2]
  \item value = [12, 10, 20, 15]
  \item capacity = 5
\end{itemize}

可以构建初步的算法如下：

\begin{lstlisting}
def knapsack(weight, value, c):
  dp = [0] * (c + 1)

  for i in range(len(weight)):
      for j in range(c, weight[i] - 1, -1):
          dp[j] = max(dp[j], dp[j - weight[i]] + value[i])

  max_value = dp[c]
\end{lstlisting}

接下来重建最优解路径。

为了重建最优解路径，我们需要在更新dp数组时记录下选择的物品。
这里可以使用一个二维数组\texttt{decision}，其中\texttt{decision[i][j]}表示在考虑前i个物品，容量为j时的选择情况。

算法如下：

\begin{lstlisting}
def knapsack(weight, value, c):
    n = len(weight)
    dp = [0] * (c + 1)
    decision = [[0] * (c + 1) for _ in range(n)]

    for i in range(n):
        for j in range(c, weight[i] - 1, -1):
            if dp[j] < dp[j - weight[i]] + value[i]:
                dp[j] = dp[j - weight[i]] + value[i]
                decision[i][j] = 1

    j = c
    chosen_items = [0] * n
    for i in range(n - 1, -1, -1):
        if decision[i][j] == 1:
            chosen_items[i] = 1
            j -= weight[i]

    # return dp[c], chosen_items
    return chosen_items
\end{lstlisting}

\subsection*{空间复杂度}
\noindent

原始的 dp 数组和 decision 数组的空间复杂度：

\begin{itemize}
  \item dp 数组的空间复杂度为 O(c)，因为我们只需要存储容量为 c 的一维数组。
  \item decision 数组的空间复杂度为 O(nc)，用于记录每个物品在每种容量下的选择情况。
\end{itemize}

由于我们只关心最优解和选择的物品，我们可以只记录选择情况，从而优化空间复杂度。

\subsection*{时间复杂度}
\noindent

对于每个物品，我们都遍历了容量从 c 到 weight[i]，因此时间复杂度为 O(nc)。

这样，我们得到了一个空间复杂度为O(nc)，
时间复杂度为O(nc)的动态规划算法来解决 0-1 背包问题，
并且能够输出最优解向量。

\end{Answer}

\begin{flushright}\begin{tabular}{rl}
  姓名: & 李俊洁\\
  班级: & 计算机2204 \\
  学号: & 20225443 \\
  课程: & 算法设计与分析 \\
  日期: & June 18, 2024 \\
  制作: & \LaTeX \\
  & 感谢您的批阅！
\end{tabular}\end{flushright}
\end{document}
